\scanpage{12}

\uline{Exemple} Soit $A$ un anneau sur $X$, soient $M,N$ respectivement les faisceaux $K^{0}$ et $K^{1}$ associés au champ additif des $A$-Modules projectifs de $\mathrm{v}.\mathrm{f}\unsure{.|\phantom{.}}$ \inserted{(p\textsuperscript{r} ex.)} Alors la construction de Mme Sinh nous fournit un champ de Picard épinglé par $M,N$, d'où une section \deleted{\unsure{du|de}} canonique du terme médian $\unsure{P|F}(M,N)$ de $(\alpha)$.\unsure{'|\textsuperscript{.}|\phantom{'}}

\uline{NB} Tout ce qui précède a les fonctorialités évidentes en $M,N,X$\ldots

\uline{Question} : Le triangle exact $(\unsure{T|J})$ et - la suite exacte $\alpha)$ sont ils connus par les compétences (Quillen, Breen, Illusie\ldots) ? \deleted{\unsure{S|X}}Connaissent-ils des variantes ``supérieures'' ? (Un principe ``géométrique'' pour en obtenir pourrait être via des $n$-champs de Picard non néc. stricts\ldots)
% archive sheet number: 11
